\begin{figure}[h]
    \centering
    \includegraphics[width=0.8\textwidth]{images/VideoPhysics.pdf}
    \caption{\small{\textbf{Illustration of poor physical commonsense by various T2V generative models.} Here, we show that the generated videos can violate a diverse range of laws pf physics such as conversation of mass, Newton's first law, and solid constitutive laws. In \name{}, we curate a wide range of prompts that would be used to assess the physical commonsense of the T2V models.}}
    \label{fig:main_examples}
\end{figure}
% \YS{for the first prompt: should it be pour milk into a cup?}

\section{Introduction}
\label{sec:introduction}

% \hb{cite VIEScore: Towards Explainable Metrics for Conditional Image Synthesis Evaluation}
% The ability to synthesize high-quality videos for a broad range of visual concepts and styles is a long-standing goal of generative modeling \cite{aldausari2022video}. In this regard, r

Recent advancements in pretraining on internet-scale video data \cite{Bain21,xue2022advancing,wang2023internvid,wang2023videofactory,chen2024panda} have led to the development of various text-to-video (T2V) generative models such as Sora \cite{liu2024sora} that can generate photo-realistic videos conditioned on a text prompt \cite{bar2024lumiere,wang2023modelscope,chen2024videocrafter2,OpenSora,singer2022make,blattmann2023align,kondratyuk2023videopoet}. Specifically, these models can generate complex scenes (e.g., `busy street in Japan') and realistic motions (e.g., `running', `pouring'), making them amenable for understanding and simulating the physical world. \hb{As humans, we develop an intuitive understanding of the object interactions through our experience with the real-world, without any formal education in physics (also termed as intuitive physics) \cite{duan2022survey}. For instance, we can predict the trajectories of the billiard balls after an application of force.} Recent efforts \cite{du2024learning,bruce2024genie} have further utilized text-guided video generation to train agents that can act, plan, and solve goals in the real world. In spite of the strong physical motivations of these works, it remains unclear \textit{how well the generated videos from T2V models adhere to the laws of physics}.


\zz{One might be tempted to assess the physical commonsense of generated videos by comparing them with physical simulations as a ground truth. However, this is non-trivial, and no similar approaches have been proposed yet. The main challenges include the lack of mature methods to accurately generate 3D geometries from single-view images or video, which is essential for physical simulations. Further, physical simulations usually require precise tuning of material parameters based on the expertise of graphics researchers to match real-world dynamics. Recently, some efforts have been made to tune simulation parameters from generated videos (e.g., \cite{huang2024dreamphysics, liu2024physics3d, zhang2024physdreamer, ni2024phyrecon}). Nevertheless, they depend on the physical plausibility of the generated videos themselves, which is again the open question we want to address. Finally, the accurate lighting and rendering are also necessary to convert physically simulated results into images and videos, yet these parameters are also unknown. Most importantly, it should be noted that physical simulations are not equivalent to ground truth. They are merely numerical solvers of differential equations that attempt to approximate and describe real-world dynamics based on models proposed by researchers.} Prior work such \hb{as VBench} \cite{huang2023vbench,liu2024evalcrafter} introduced a comprehensive benchmark to evaluate various qualities of generated videos (e.g., motion smoothness, background consistency) using existing models, but it does not specifically address the generated videos' adherence to physical laws. Therefore, existing benchmarks and metrics are either unreliable or lack coverage for holistic evaluation of the physical commonsense capabilities.



To this end, we propose \name{}, a dataset designed to evaluate the adherence of generated videos to physical commonsense in real-world scenarios. Specifically, we focus on the intuitive understanding of the behavior and dynamics of various states of matter (solids, fluids) in the physical world \cite{piloto2022intuitive,yildirim2024perception,bisk2020piqa}. For instance, `water pouring into a glass' will intuitively result in the water level in the glass rising over time. As a result, we rely on human perception and experience in the physical world to assess the adherence of the generated videos to physical laws instead of precise dynamical equations, which are harder to assess. In Figure \ref{fig:main_examples}, we provide qualitative examples to illustrate physical commonsense violations in the videos. Our dataset is constructed through a three-stage pipeline that involves (a) prompting a large language model \cite{gpt4} to generate candidate captions that depict interactions between diverse states of matter (e.g., solid-solid, solid-fluid, fluid-fluid), (b) human verification of the generated captions, and (c) annotating the complexity in rendering objects or synthesizing motions described in the captions based on physics simulation. 


In total, \name{} comprises 688 high-quality, human-verified captions that will be used to generate videos from T2V models. In addition, the dataset consists human-labeled annotations for physical commonsense of the generated videos. Specifically, we acquire generated videos from \hb{\textbf{twelve}} diverse T2V models including open models (e.g., OpenSora \cite{OpenSora}, SVD \cite{blattmann2023stable}, VideoCrafter2 \cite{chen2024videocrafter2}, CogVideoX \cite{yang2024cogvideox}) and closed models (e.g., Pika \cite{pika}, Lumiere \cite{bar2024lumiere}, Gen-2 \cite{gen2}, \hb{Dream Machine \cite{lumalabsLumaDream}}). Subsequently, we perform human evaluation on the generated videos for semantic adherence to the conditioning text (e.g., do the videos follow the caption?) and physical commonsense (e.g., do the videos follow physical laws intuitively?). Interestingly, we find that the existing T2V generative models severely lack the capability to follow caption accurately and generate videos with physical commonsense. Specifically, the best performing model, \hb{CogVideoX-5B}, follows the text and generates physically accurate videos for $39.6\%$ of the instances (\S \ref{sec:results}). \hb{Our fine-grained analysis reveals that current T2V models are particularly poor at generating physically plausible videos for prompts that require solid-solid interaction (e.g., ball bouncing on the floor, hammer hits a nail).}
In Figure \ref{fig:teaser_graph}, we compare the performance (i.e., accurate semantic adherence and physical commonsense) of various T2V generative models on the \name{} dataset. \hb{In addition, we perform a detailed qualitative analysis to study the modes of the failures for different models in detail (\S \ref{sec:qualitative_examples})}. \tx{In particular, we observe that the models often struggle to accurately identify individual objects and comprehend their material properties, which is essential for generating physically plausible dynamics. For instance, an object recognized as a rigid body in the physical world should not deform over time.}

Although human evaluation of semantic adherence and physical commonsense is reliable, it is both expensive and difficult to scale. To address this challenge, we introduce \model{}, an \hb{open} video-language model designed to assess the semantic adherence and physical commonsense of generated videos using user queries grounded in text (\S \ref{sec:automatic_evaluation_results}). Specifically, we finetune \videocon{} \cite{bansal2023videocon}, a robust semantic adherence evaluator for real videos, on generated videos and human annotations \hb{collected as a part of our dataset}. Our results demonstrate that \model{} outperforms Gemini-Pro-Vision-1.5 \cite{reid2024gemini}, showing a $9$ points improvement in semantic adherence and a $15$ points improvement in physical commonsense on unseen prompts. \hb{Further, we show that \model{} generalizes to unseen generative models, which established its reliability for evaluating future generative models.} Overall, the \name{} dataset aims to bridge the gap in understanding physical commonsense in generated videos and enables scalable testing \hb{for upcoming T2V models}.